% ChunkTest.tex — regression test for writeChunk / showChunk / runRChunk
%
% Covers three bugs fixed in v2.4:
%   1. Stream exhaustion: \newread/\newwrite inside writeChunk exhausted
%      TeX's 16-stream limit. Fixed by pre-allocating streams at package level.
%      Tested here by calling \showChunk 18 times (> the old limit of 16).
%   2. Unknown shell command: showChunk called \write18{ad ...} which is not
%      a standard command. Fixed by removing the call (writeChunk handles
%      directory creation itself now).
%   3. Subdirectory path: writeChunk wrote chunk files to generated/<path>
%      without creating the intermediate directory. Fixed by running
%      mkdir -p on the directory component of the source path.
%      Tested here by using a source file in the code/ subdirectory.

\documentclass[12pt]{article}
\usepackage[nohup,R,stopserver,listings]{../../runcode}
\usepackage[margin=1in]{geometry}

\begin{document}

\section*{Test 1: \texttt{showChunk} with a file in a subdirectory}

Source file is \texttt{code/stats.R}.
The chunk output must land in \texttt{generated/code/stats.R-chunkXX.txt},
requiring \texttt{generated/code/} to be created automatically.

\showChunk{R}{code/stats.R}{chunk01}
\showChunk{R}{code/stats.R}{chunk02}
\showChunk{R}{code/stats.R}{chunk03}

\section*{Test 2: \texttt{showChunk} with a flat file (no subdirectory)}

Source file is \texttt{flatcode.R} in the current directory.
Output lands in \texttt{generated/flatcode.R-flat\_hello.txt}.

\showChunk{R}{flatcode.R}{flat_hello}
\showChunk{R}{flatcode.R}{flat_arithmetic}

\section*{Test 3: Stream exhaustion — 18 \texttt{showChunk} calls}

Calling \texttt{showChunk} 18 times on the same file exercises the
pre-allocated stream fix. The old code would crash after the 16th call
with ``No room for a new read stream.''

\showChunk{R}{code/stats.R}{chunk04}
\showChunk{R}{code/stats.R}{chunk05}
\showChunk{R}{code/stats.R}{chunk06}
\showChunk{R}{code/stats.R}{chunk07}
\showChunk{R}{code/stats.R}{chunk08}
\showChunk{R}{code/stats.R}{chunk09}
\showChunk{R}{code/stats.R}{chunk10}
\showChunk{R}{code/stats.R}{chunk11}
\showChunk{R}{code/stats.R}{chunk12}
\showChunk{R}{code/stats.R}{chunk13}
\showChunk{R}{code/stats.R}{chunk14}
\showChunk{R}{code/stats.R}{chunk15}
\showChunk{R}{code/stats.R}{chunk16}
\showChunk{R}{code/stats.R}{chunk17}
\showChunk{R}{code/stats.R}{chunk18}

\section*{Test 4: \texttt{runRChunk} with a subdirectory file}

\texttt{runRChunk} should extract the chunk, run it via the R server,
and display the output.

\runRChunk{code/stats.R}{chunk01}
\runRChunk{code/stats.R}{chunk02}
\runRChunk{code/stats.R}{chunk17}

\section*{Test 5: \texttt{runRChunk} with a flat file}

\runRChunk{flatcode.R}{flat_hello}
\runRChunk{flatcode.R}{flat_arithmetic}

\section*{Test 6: Missing chunk label (should show red error, not crash)}

\showChunk{R}{code/stats.R}{no_such_label}

\section*{Test 7: Missing source file (should show red error, not crash)}

\showChunk{R}{no_such_file.R}{chunk01}

\end{document}
